% =====================================================================
% PASTE FILE -- abstract nou, scris pentru introducerea de la revizia 2.
% Se copiaza peste blocul \begin{abstract}...\end{abstract} din
% acl_latex.tex in Overleaf.
%
% Cel vechi (report/report.tex) era de raport de curs: descria ce s-a
% facut, fara niciun rezultat. Asta il inlocuieste. ~140 de cuvinte.
% Anonimizat: fara nume, fara link.
% =====================================================================

\begin{abstract}
Lexical complexity prediction is normally solved by fine-tuning an
encoder on annotated data in the target language, which excludes every
language that has none. We predict the lexical complexity of Romanian
words in context without using a single Romanian training label, by
fusing four signals that carry complexity knowledge from outside the
language: cross-lingual transfer through a frozen multilingual encoder
trained on English and Romance annotations, a complexity rating and
token surprisal from a Romanian language model, and a reading-effort
score derived from eye-tracking. The system reaches $69.74$ on the
shared-task metric, against $10.27$ for surface features alone, and
Pearson $r = 0.913$ on held-out Romanian items. The largest single gain
($+4.09$) comes from a lemma-frequency correction that repairs the
frequency estimates inflection corrupts, and transfers to any
morphologically rich language served by the same frequency lists. We
also report two negative results: with no target labels, fine-tuning
the encoder is consistently worse than freezing it, and in-source
validation does not predict transfer quality.
\end{abstract}
